\section{复算与人工智能使用说明}
支撑材料给出三问的完整调度文本。问题一只含原始节点顺序；问题二、三分别提供扩展节点顺序、初始地址偏移与按换出顺序列出的新偏移。新增节点每对由换出及换入组成，重新排列后也保持列表与节点编号一致。复算应从原图和这些文本重新建立依赖、空间状态及执行时间，不能只读取汇总数值。

程序以 Python 实现。图读取、顺序生成、地址分配、流水重排与附件重放分模块组织；候选规则及平局处理均为确定性。复算材料保留全部候选与失败记录，论文表中的最终值来自同一组六图方案。新增过程分析从这些方案重放得到，并以九节点构造例的全部 126 个拓扑序检验生命周期峰值。

研究过程中使用 OpenAI Codex 辅助模型推导、程序实现、结果分析及论文编辑。历史求解回合的精确模型标识和版本发布日期未核实，参赛队员独立复核尚未完成。若用于参赛，队员需依据适用规定理解并核验相关内容，以自身语言形成最终文本，并补齐真实工具记录与使用声明\cite{ai}。
